\documentclass[
    12pt,
    oneside,
    a4paper,
    brazil
]{abntex2}

\usepackage[utf8]{inputenc}
\usepackage{mathptmx}
\usepackage[T1]{fontenc}
\usepackage{indentfirst}
\usepackage{setspace}

% Informações
\titulo{Documento de Requisitos}
\autor{Caio Collino, Thiago Lins, Vinicius Vianna}
\local{São Paulo}
\data{2026}

\instituicao{
  Centro Universitário Senac
  \par
  Ciência da Computação
}

\hypersetup{
    hidelinks
}

\tipotrabalho{Documento Técnico}
\preambulo{Documento de requisitos funcionais e não funcionais desenvolvido para a disciplina de Engenharia de Software.}

\begin{document}

\imprimircapa

\tableofcontents*
\cleardoublepage

\OnehalfSpacing

\chapter{Introdução}
Este documento visa listar os requisitos funcionais e não funcionais da plataforma de reserva de hotéis \textbf{StayHub}, um sistema que centralizará opções de estadia e seus preços, de diferentes sites.

\chapter{Requisitos Funcionais}
\section{RF01 - Exemplo}
Descrição do requisito funcional.

\chapter{Requisitos Não Funcionais}
\section{RNF01 - Desempenho}
Descrição do requisito não funcional.

\chapter{Conclusão}

\end{document}
